\section*{अध्याय \ref{ch:b3:asymmetric-synthesis} --- असममित संश्लेषण}

\begin{solution}{exo:b3:asymmetric-synthesis:1}
ee 80\,\%: $\mathrm{er} = 1.8/0.2 = 9$ (90\,:\,10)। er 99\,:\,1: $\mathrm{ee} = 98/100 = 98$\,\%। ee 50\,\%: 75\,\% और 25\,\%।
\end{solution}

\begin{solution}{exo:b3:asymmetric-synthesis:2}
$(1840 - 60)/(1840 + 60) = 0.937$: ee 93.7\,\%, $1840/60 = 31$ का er।
\end{solution}

\begin{solution}{exo:b3:asymmetric-synthesis:3}
प्राथमिकताएँ O $>$ एथिल $>$ H। O ऊपर, एथिल नीचे बाएँ और H नीचे दाएँ बनाने पर,
सामने का फलक \textit{Si} है, पीछे का \textit{Re}। \omterm{def:b3:asymmetric-synthesis:faces}{Re फलक} से जोड़ा गया मेथिल समूह,
उत्पाद में प्राथमिकताओं OH $>$ एथिल $>$ मेथिल $>$ H के साथ, (\textit{R})-ब्यूटेन-2-ऑल देता है।
\end{solution}

\begin{solution}{exo:b3:asymmetric-synthesis:4}
एथेनॉल: \omterm{def:b3:asymmetric-synthesis:topicity}{प्रतिबिंबस्थ} (अणु के \omterm{def:b3:point-groups:mirror}{दर्पण तल} द्वारा परस्पर बदले जाते हैं; अकिरैल विलायक में
समान विस्थापन)। ब्यूटेन-2-ऑल: \omterm{def:b3:asymmetric-synthesis:topicity}{अप्रतिबिंबस्थ} (त्रिविमजनक केंद्र के बगल में;
भिन्न विस्थापन)।
\end{solution}

\begin{solution}{exo:b3:asymmetric-synthesis:5}
er 199: \qty{298}{K} पर $RT\ln 199 = \qty{13.1}{kJ/mol}$, \qty{195}{K} पर \qty{8.6}{kJ/mol}। \qty{25}{\celsius} पर 80\,\% ee, er 9 है,
$\Delta\Delta G^\ddagger = RT\ln 9 = \qty{5.45}{kJ/mol}$; \qty{-78}{\celsius} पर $\mathrm{er} = \eu^{5450/(8.314 \times 195.15)} = 28.7$, ee 93\,\%।
\end{solution}

\begin{solution}{exo:b3:asymmetric-synthesis:6}
\omterm{def:b3:asymmetric-synthesis:ee}{प्रकाशिक शुद्धता} $12.0/15.0 = 80$\,\%: 90\,\% (\textit{S}) और 10\,\% (\textit{R})। छोटा घूर्णन, अशुद्धियाँ,
सांद्रता और विलायक प्रभाव, और घूर्णन--संघटन का अरैखिक
व्यवहार इसे अविश्वसनीय बनाते हैं; वर्णलेखन प्रतिबिंबरूपियों को सीधे पृथक करके
गिनता है।
\end{solution}

\begin{solution}{exo:b3:asymmetric-synthesis:7}
$s = \ln(0.40 \times 0.10)/\ln(0.40 \times 1.90) = -3.22/-0.274 = 12$।
\end{solution}

\begin{solution}{exo:b3:asymmetric-synthesis:8}\emergencystretch=3em
L-($+$)-डाइएथिल टार्ट्रेट के साथ (2\textit{S},3\textit{S})-2,3-एपॉक्सीहेक्सेन-1-ऑल; D-($-$)-डाइएथिल टार्ट्रेट के साथ
प्रतिबिंबरूपी, (2\textit{R},3\textit{R})।
\end{solution}

\begin{solution}{exo:b3:asymmetric-synthesis:9}
L = फ़ेनिल, M = मेथिल, S = H। (\textit{S}) ऐल्डिहाइड के लिए क्रियाशील संरूपी में फ़ेनिल
कार्बोनिल तल के लंबवत होता है और हाइड्रोजन नाभिकरागी के पथ के बगल में,
जो फ़ेनिल के anti आक्रमण करता है। विन्यासों को निकालने पर मुख्य अप्रतिबिंबरूपी
(2\textit{S},3\textit{S})-3-फ़ेनिलब्यूटेन-2-ऑल मिलता है।
\end{solution}

\begin{solution}{exo:b3:asymmetric-synthesis:10}
आरंभ में दोनों प्रतिबिंबरूपी समान मात्राओं में होते हैं, अतः उत्पाद
$k_{\mathrm{fast}} : k_{\mathrm{slow}} = s$ अनुपात में बनते हैं: $\mathrm{ee}_p = (s - 1)/(s + 1)$, $s = 20$ के लिए 0.905। अभिक्रिया आगे बढ़ने पर
तीव्र प्रतिबिंबरूपी घटता जाता है और उत्पाद का ee गिरता है, जबकि क्रियाधार का ee 1 की ओर
बढ़ता है: क्रियाधार को आगे जाकर सदैव शुद्ध किया जा सकता है; उत्पाद को
नहीं।
\end{solution}

\begin{solution}{exo:b3:asymmetric-synthesis:11}
तीव्र रेसेमीकरण के साथ: 100\,\% तक लब्धि, $\mathrm{er} = 99$, ee 98\,\%। इसके बिना, 50\,\%
रूपांतरण पर: लब्धि अधिक से अधिक 50\,\%, और उत्पाद ee 93\,\% (क्रियाधार का ee भी
वही, 93\,\%)।
\end{solution}

\begin{solution}{exo:b3:asymmetric-synthesis:12}
जब दोनों संकुल अपनी अभिक्रिया से तेज़ी से परस्पर रूपांतरित होते हैं, तो उत्पाद अनुपात
$k_B[\mathrm B]/(k_A[\mathrm A]) = 600/10 = 60$ है: गौण संकुल B उत्पाद का 98\,\% देता है (ee 97\,\%)।
वरणात्मकता दोनों संक्रमण अवस्थाओं के अंतर से तय होती है, न कि
मध्यवर्तियों की \omterm{def:b3:partition-functions:boltzmann}{समष्टियों} से।
\end{solution}

\begin{solution}{pb:b3:asymmetric-synthesis:1}
\textbf{1.} नाइट्रोजन-धारी ऐल्कीन कार्बन पर तीन भिन्न प्रतिस्थापी
और दो भिन्न फलक हैं; उस पर जुड़ा हाइड्रोजन ऐमीनो
अम्ल के $\alpha$-कार्बन को त्रिविमजनक केंद्र बना देता है।

\textbf{2.} रेसेमिक मिश्रण: दोनों फलकों पर समान दर से आक्रमण होता है।

\textbf{3.} इसका कार्बोनिल ऑक्सीजन C=C आबंध के साथ रोडियम से बँधता है: क्रियाधार
कीलेट के रूप में, एक नियत ज्यामिति में, थामा जाता है जिसमें लिगैंड विभेद कर सकता है।

\textbf{4.} यह धातु के चारों ओर एक किरैल जेब बनाता है ($C_2$-सममित): क्रियाधार को
एक या दूसरे फलक से बाँधने के दोनों ढंग अप्रतिबिंबरूपी हैं, भिन्न
ऊर्जाओं और अभिक्रियाशीलताओं के साथ।

\textbf{5.} नहीं: गौण संकुल हाइड्रोजन से कहीं तेज़ी से अभिक्रिया करता है (कर्टिन--हैमेट)।

\textbf{6.} धातु बँधे हुए ऐल्कीन के एक फलक पर होती है, और हाइड्रोजन
धातु से स्थानांतरित होते हैं।

\textbf{7.} $\mathrm{er} = 1.95/0.05 = 39$।

\textbf{8.} $\Delta\Delta G^\ddagger = RT\ln 39 = 8.314 \times 298.15 \times 3.664 = \qty{9.1}{kJ/mol}$।

\textbf{9.} $\mathrm{er} = \eu^{9082/(8.314 \times 195.15)} = 270$: ee 99.3\,\%।

\textbf{10.} $\mathrm{er} = \eu^{11082/(8.314 \times 298.15)} = 87$: ee 97.7\,\%।

\textbf{11.} एक विशिष्ट हाइड्रोजन आबंध से कम: त्रिविम संपर्कों में एक छोटा अंतर
या क्रियाधार और लिगैंड के बीच एक दुर्बल अन्योन्यक्रिया परिणाम तय करती है,
इसीलिए ऐसे उत्प्रेरकों का मूल सिद्धांतों से अभिकल्पन कठिन है।

\textbf{12.} दर गिरती है, हाइड्रोजन का विलयन और उत्प्रेरक का व्यवहार
निम्न ताप पर बदल जाते हैं, और किसी बड़े रिएक्टर को ठंडा करने में ऊर्जा लगती है।

\textbf{13.} 97.5 g (\textit{S}) और 2.5 g (\textit{R}) प्रतिबिंबरूपी।

\textbf{14.} $97.5 - 5.0 = \qty{92.5}{g}$ क्रिस्टल, व्यावहारिक रूप से प्रतिबिंबरूपी-शुद्ध।

\textbf{15.} अपरिष्कृत का $92.5/100 = 92.5$\,\% ((\textit{S}) प्रतिबिंबरूपी का 95\,\%)।

\textbf{16.} क्रिस्टल आधिक्य वाले प्रतिबिंबरूपी को लेते हैं, रेसेमिक भाग गौण प्रतिबिंबरूपी के साथ
विलयन में रहता है; रेसेमिक मिश्रण में एकत्र करने के लिए कोई आधिक्य नहीं होता, और
क्रिस्टलन रेसेमिक क्रिस्टल देता है।

\textbf{17.} रेसेमिक संदर्भ के विरुद्ध किरैल HPLC द्वारा।

\textbf{18.} 50\,\%: केवल धीमा प्रतिबिंबरूपी बचता है।

\textbf{19.} $s = \ln(0.45 \times 0.05)/\ln(0.45 \times 1.95) = -3.79/-0.131 = 29$।

\textbf{20.} रेसेमिक मिश्रण का 45\,\% (आरंभ में उपस्थित वांछित प्रतिबिंबरूपी का 88\,\%,
95\,\% ee पर)।

\textbf{21.} क्रियाधार के तीव्र रेसेमीकरण के साथ लब्धि 100\,\% के निकट पहुँच सकती है।

\textbf{22.} यह सस्ते अकिरैल पूर्वगामी के सारे भाग को सही प्रतिबिंबरूपी में बदल देता है,
अत्यल्प मात्रा में प्रयुक्त उत्प्रेरक के साथ और कोई प्रतिबिंबरूपी फेंके बिना।

\textbf{23.} रोडियम विषैला है और औषधों में इसके अवशेष प्रति दस लाख कुछ भागों तक
सीमित हैं; इन्हें प्रेरणिक युग्मित प्लाज़्मा द्रव्यमान स्पेक्ट्रममिति से मापा जाता है।

\textbf{24.} \qty{25}{\celsius} पर 95\,\% ee, $\Delta\Delta G^\ddagger = RT\ln 39 \approx \qty{9.1}{kJ/mol}$ के संगत है।
\end{solution}
